\documentclass[10pt]{article}
\usepackage[a4paper,margin=1.7cm]{geometry}
\usepackage[T1]{fontenc}
\usepackage{lmodern}
\usepackage{amsmath,amssymb}
\usepackage{longtable,booktabs,array}
\usepackage{listings}
\usepackage[hidelinks]{hyperref}
\lstset{basicstyle=\ttfamily\footnotesize,breaklines=true,columns=fullflexible,keepspaces=true,upquote=true}
\makeatletter
\newcommand\keepnext{\par\nopagebreak\@nobreaktrue\everypar{\@nobreakfalse\everypar{}}}
\newcommand\mdneed[1]{\par\if@nobreak\else\begingroup\setlength{\dimen@}{#1}\dimen@ii\pagegoal\advance\dimen@ii-\pagetotal
  \ifdim\dimen@>\dimen@ii\ifdim\dimen@ii>\z@\vfil\fi\break\fi\endgroup\fi}
\def\md@ltpatch#1\penalty\z@#2\md@ltend{\def\LT@start{#1\if@nobreak\penalty\@M\else\penalty\z@\fi#2}}
\expandafter\md@ltpatch\LT@start\md@ltend
\makeatother
\setcounter{MaxMatrixCols}{20}
\allowdisplaybreaks
\setlength{\parindent}{0pt}\setlength{\parskip}{4pt}\setlength{\emergencystretch}{3em}\sloppy\raggedright
\begin{document}

{\raggedright\LARGE\bfseries Provenance of `Einstein-Rosen bridges in the gravitational fields of dirac16complex'\par}\medskip

\section*{1. Sources}

\begin{itemize}
\item A. Einstein and N. Rosen, ``The Particle Problem in the General Theory of Relativity'', Phys. Rev. 48, 73-77 (1935), doi:10.1103/PhysRev.48.73. Quoted in Pre-gravity6plane.nb (cells 1397-1399, 1406, 1412): the substitution $u^2 = r - 2m$, two congruent sheets $u>0$, $u<0$, joined by the hyperplane $u = 0$ in which $g$ vanishes, called a bridge. The charged bridge with reversed sign of the electromagnetic energy and $m=0$ is from the same paper; here it is re-derived as $r^2 = u^2 + q^2$, $ds^2 = -\frac{u^2}{u^2+q^2}dt^2 + du^2 + (u^2+q^2)d\Omega^2$ and verified.
\item Pre-gravity6plane.nb (P. L. Nash, (C) 2022, GPL): primordial vielbein $\mathrm{diag}(\tan 6Hx_0,\ e^{-a_4}\sin^{-1/6}6Hx_0\ (\times3),\ 1,\ e^{a_4}\sin^{-1/6}6Hx_0\ (\times3))$ (cells 901, 913) and the domain $0<6Hx_0<\pi/2$ (cell 48).
\item The modern and future fields are this work's definitions ($f = 1 - r_s/x_0$, $f = 1 - r_s/x_0 - H_\Lambda^2x_0^2$ for the $(x_0,x_4)$ metric $dx_0^2/f - f\,dx_4^2$).
\item Standard results used in Section 4: W. G. Unruh, Phys. Rev. D 14, 870 (1976); W. Israel, Phys. Lett. A 57, 107 (1976); J. B. Hartle and S. W. Hawking, Phys. Rev. D 13, 2188 (1976); Kruskal coordinates (M. D. Kruskal, Phys. Rev. 119, 1743 (1960)). The picture of bridges as pair correlations follows J. Maldacena and L. Susskind, Fortsch. Phys. 61, 781 (2013).
\end{itemize}

\section*{2. Computations (Wolfram Language 15.0.1, exact, 2026-10-30)}

{\small\begin{longtable}{>{\raggedright\arraybackslash}p{\dimexpr 0.7323\linewidth-2\tabcolsep\relax}>{\raggedright\arraybackslash}p{\dimexpr 0.2277\linewidth-2\tabcolsep\relax}}\toprule
result & 6planes.nb \\\midrule\endhead
\bottomrule\endfoot
neutral bridge: $R_{\mu\nu}=0$, $\det g = -4u^2(u^2+2m)^4\sin^2\theta$, even in $u$, Kretschmann $48m^2/(u^2+2m)^6$ & Part 8, Out[31] \\*
charged bridge: $G_{\mu\nu}+8\pi T^{EM}_{\mu\nu}=0$, $\det g(0)=0$, Kretschmann$(0) = 56/q^4$ & Part 8, Out[31] \\*
modern and future bridges: $g_{uu}(0)$, $g_{44}(0)$, evenness, $R(0)$, $\kappa$ (exact cubic roots; check at $u=10^{-6}$) & Part 8.1, Out[32] \\*
primordial: $R = -\frac{3H^2}{2}(57+64\cos12Hx_0+3\cos24Hx_0)\csc^46Hx_0$, series $-\frac{31}{216H^2x_0^4}$, $\det g = \sec^2 6Hx_0$ & Part 8.2, Out[33] \\*
continuation factor $e^{-\pi\epsilon/\kappa}$, occupation $1/(e^{2\pi\epsilon/\kappa}+1)$ equals Fermi-Dirac at $T_H = \kappa/2\pi$ & Part 8.3, Out[34] \\*
\end{longtable}}

The curvature routines (Christoffel symbols, Riemann tensor $R^i{}_{jkl} = \partial_k\Gamma^i{}_{jl} - \partial_l\Gamma^i{}_{jk} + \Gamma^i{}_{km}\Gamma^m{}_{jl} - \Gamma^i{}_{lm}\Gamma^m{}_{jk}$, Ricci tensor $R_{jl} = R^i{}_{jil}$, Einstein tensor, Kretschmann scalar) are written out in 6planes.nb and were first tested on the Einstein-Rosen metrics in a scratch script (test\_er.wls) with identical results.


\section*{3. New in this work}

Application of the Einstein-Rosen substitution $x_0 = x_h + su^2$ to the modern and future fields; the proof that the primordial field's degenerate hypersurface is singular; the identification of the Kruskal pairing with the $\{+\epsilon,-\epsilon\}$ universes of dirac16complex.


\end{document}